% ============================================
% DOC-18 (EN)
% English translation — SPEC-935-R355
% ============================================
\section*{\textbf{DOC-18} --- Official Report on the Pen}
\markboth{DOC-18 --- Official Report on the Pen}{}

\noindent \textit{1916 --- Reconstruction concerning the Alagadiço Camp}

\noindent \textbf{Fictional reconstruction --- no counterpart located in the National Archive}
\textbf{1916 --- Narrative document, not a transcription}

\vspace{0.5em}\noindent\rule{\textwidth}{0.5pt}\vspace{0.5em}

\noindent \textit{FICTIONAL RECONSTRUCTION OF AN ADMINISTRATIVE RECORD}
\textit{PUBLIC RELIEF}
\textit{NARRATIVE REPORT ON THE CAMPS OF CEARÁ}
\textit{FISCAL YEAR 1915-1916}

\vspace{0.5em}\noindent\rule{\textwidth}{0.5pt}\vspace{0.5em}

\noindent \textbf{Camp:} Alagadiço, near Fortaleza
\textbf{Official capacity:} not stated in this reconstruction
\textbf{Number interned during the period:} imprecise; the sources consulted record thousands of confined people
\textbf{Deaths:} no verifiable total reproduced in this narrative

\noindent \textbf{Voice of the fictional reporting officer:}

\noindent \textit{\textquotedblleft{}The concentration camps, officially called 'government pens', were established as an emergency measure to contain the flow of retirantes toward the capital. In practice, they functioned as warehouses of human beings in a state of starvation.}

\noindent \textit{Food and sanitary conditions were insufficient to protect the lives of those confined. This reconstruction attributes neither a mortality percentage nor documentary authorship to this wording.}

\noindent \textit{The closure of the camps and the adoption of effective relief measures are recommended.\textquotedblright{}}

\vspace{0.5em}\noindent\rule{\textwidth}{0.5pt}\vspace{0.5em}

\noindent \textbf{Editorial note on the reconstruction:}

\noindent The existence of the Alagadiço Camp, near Fortaleza, during the 1915 drought is documented. The numbers, the reporting voice, and the administrative form above are fictional devices. This text does not claim that the camp remained open until 1917; Joaquim's transfer that year belongs to the character's fictional chronology.

\vspace{0.5em}
\noindent\textit{\textbf{↪ Routes:} \rota{DOC-19} \rotasep \rota{MEM-06}}
